\documentclass[11pt,a4paper]{article}
\usepackage[T1]{fontenc}
\usepackage[utf8]{inputenc}
\usepackage{amsmath,amssymb,amsthm}
\usepackage[margin=1in]{geometry}
\usepackage[colorlinks=true,linkcolor=blue,urlcolor=blue]{hyperref}
\theoremstyle{plain}
\newtheorem{theorem}{Theorem}
\newtheorem{lemma}[theorem]{Lemma}
\newtheorem{proposition}[theorem]{Proposition}
\newtheorem{corollary}[theorem]{Corollary}
\theoremstyle{definition}
\newtheorem*{remark}{Remark}

\title{One local factor:\\
a proof of Mathar's M\"obius-transform conjecture on\\
OEIS A062368}
\author{Adrian Perez Fontelles\\ \small Independent researcher}
\date{24 August 2026}
\begin{document}
\maketitle

\begin{abstract}
OEIS A062368 is the multiplicative function with $a(p^{e})=(e+1)(e+2)(4e+3)/6$. In August
2012 R.~J. Mathar conjectured that it is the third inverse M\"obius transform of
$4^{\omega(n)}$, where $\omega$ counts distinct prime factors. We prove it. Both sides are
multiplicative, so the whole statement is a single identity between local Dirichlet
factors, and that identity is $\binom{e+3}{3}+3\binom{e+2}{3}=(e+1)(e+2)(4e+3)/6$: the
generating function of $a(p^{e})$ is $(1+3x)/(1-x)^{4}$, which is exactly the local factor
of $4^{\omega}$ times three copies of $\zeta$.
\end{abstract}

\section{The sequence and the conjecture}

OEIS A062368 (Vladeta Jovovic, 7 July 2001) has offset 1 and is defined as the
multiplicative function with
\begin{equation}\label{eq:def}
  a(p^{e})=\frac{(e+1)(e+2)(4e+3)}{6},
\end{equation}
so $a=1,7,7,22,7,49,7,50,22,49,\dots$. The entry carries

\begin{quote}
\textbf{Conjecture:} this is the third inverse M\"obius transform of the sequence
$4^{\mathrm{A001221}(n)}$. --- R. J. Mathar, Aug 09 2012
\end{quote}

Here $\mathrm{A001221}(n)=\omega(n)$ is the number of distinct primes dividing $n$, and the
inverse M\"obius transform of an arithmetic function $b$ is $(b*1)(n)=\sum_{d\mid n}b(d)$,
where $*$ is Dirichlet convolution and $1$ is the constant function 1. Applying it three
times gives $b*1*1*1=b*\tau_{3}$, with $\tau_{3}=1*1*1=\mathrm{A007425}$. So the conjecture
reads
\begin{equation}\label{eq:conj}
  a=4^{\omega}*1*1*1,\qquad\text{that is}\qquad
  a(n)=\sum_{d\mid n}4^{\omega(d)}\,\tau_{3}(n/d).
\end{equation}
As of the ``Last modified'' line on the live entry (24 August 2026) the statement is still
recorded as an unproven conjecture, with no proof in the entry's links or programs.

\begin{remark}[status, stated plainly]
On 10 July 2026 R.~Oudra added to the FORMULA section of the same entry the line
$a(n)=\sum_{d\mid n}4^{\omega(d)}\tau_{3}(n/d)$, unlabelled and without proof, together
with the Dirichlet generating function $\zeta(s)^{4}\prod_{p}(1+3p^{-s})$. That formula is
\eqref{eq:conj}, that is Mathar's conjecture restated; it was added six weeks before this
note and the conjecture label was not removed. I have found no proof of either on the entry
or in the literature. What follows supplies one, and shows the two statements are the same.
\end{remark}

\section{Proof}

Write $x=p^{-s}$ and, for an arithmetic function $b$, let $B_{p}(x)=\sum_{e\ge0}b(p^{e})
x^{e}$ denote its local factor at $p$; for multiplicative $b$ the Dirichlet series is
$\sum_{n}b(n)n^{-s}=\prod_{p}B_{p}(p^{-s})$, and Dirichlet convolution corresponds to
multiplication of local factors. All series below are formal.

\begin{lemma}\label{lem:four}
The local factor of $4^{\omega}$ at every prime is $(1+3x)/(1-x)$, and hence the local
factor of $4^{\omega}*1*1*1$ is $(1+3x)/(1-x)^{4}$.
\end{lemma}

\begin{proof}
$\omega(p^{e})=1$ for $e\ge1$ and $\omega(1)=0$, so
\[
  \sum_{e\ge0}4^{\omega(p^{e})}x^{e}=1+4\sum_{e\ge1}x^{e}=1+\frac{4x}{1-x}=\frac{1+3x}{1-x}.
\]
The local factor of $1$ is $\sum_{e\ge0}x^{e}=1/(1-x)$; multiplying by it three times gives
the second claim.
\end{proof}

\begin{lemma}\label{lem:a}
The local factor of $a$ at every prime is $(1+3x)/(1-x)^{4}$.
\end{lemma}

\begin{proof}
Since $(1-x)^{-4}=\sum_{e\ge0}\binom{e+3}{3}x^{e}$,
\[
  \frac{1+3x}{(1-x)^{4}}=\sum_{e\ge0}\left[\binom{e+3}{3}+3\binom{e+2}{3}\right]x^{e},
\]
with the convention $\binom{2}{3}=0$ at $e=0$. Now
\[
  \binom{e+3}{3}+3\binom{e+2}{3}=\frac{(e+1)(e+2)(e+3)}{6}+3\cdot\frac{e(e+1)(e+2)}{6}
  =\frac{(e+1)(e+2)\bigl[(e+3)+3e\bigr]}{6}=\frac{(e+1)(e+2)(4e+3)}{6},
\]
which is $a(p^{e})$ by \eqref{eq:def}.
\end{proof}

\begin{theorem}\label{thm:main}
$a=4^{\omega}*1*1*1$; that is, Mathar's conjecture holds:
\[
  a(n)=\sum_{d\mid n}4^{\omega(d)}\,\tau_{3}(n/d)\qquad\text{for every }n\ge1.
\]
\end{theorem}

\begin{proof}
Both sides are multiplicative --- $a$ by definition, and the right-hand side because
Dirichlet convolution of multiplicative functions is multiplicative. Two multiplicative
functions are equal precisely when they agree at every prime power, equivalently when their
local factors agree at every prime. By Lemmas~\ref{lem:four} and \ref{lem:a} both local
factors equal $(1+3x)/(1-x)^{4}$ at every prime.
\end{proof}

\begin{corollary}
The Dirichlet generating function of A062368 is $\zeta(s)^{4}\prod_{p}\bigl(1+3p^{-s}\bigr)$.
\end{corollary}

\begin{proof}
Take the product of the local factors of Lemma~\ref{lem:a} over all primes and use
$\prod_{p}(1-p^{-s})^{-1}=\zeta(s)$.
\end{proof}

\begin{remark}[why the number three appears twice, and independently]
The conjecture involves three copies of $1$; the constant 3 in the local factor $(1+3x)$
comes from $4-1$, not from the number of transforms. They are unrelated, which is
presumably why the identity looked opaque enough to be left as a conjecture: the statement
pairs the multiplier 4 in $4^{\omega}$ with the exponent 4 in $(1-x)^{-4}$, and those two
fours have different origins --- one is $1+3$, the other is $1+3$ copies of $\zeta$.
\end{remark}

\begin{remark}[honest assessment]
Half a page, and nothing beyond multiplicativity and one binomial identity. The natural
destination is a comment on the entry. It is worth recording because the statement has
carried the label ``Conjecture'' since 2012, and because the same identity was added to the
same entry's FORMULA section as an unlabelled assertion in July 2026 without the conjecture
being marked resolved --- the same pattern seen on A005329 and A129364, where a later
contributor restates a standing conjecture as a fact and nobody reconciles the two.
\end{remark}

\section{Verification}

A verification script, written separately from this note, checks every claim and prints
every number quoted. It (i) reproduces the first 68 published terms of A062368 from the
multiplicative definition \eqref{eq:def}; (ii) confirms the conjecture by applying the
inverse M\"obius transform to $4^{\omega}$ three times, literally, and comparing with
$a(n)$ for $1\le n\le300$, with no discrepancy, and likewise in the convolved form
$\sum_{d\mid n}4^{\omega(d)}\tau_{3}(n/d)$; (iii) confirms the local-factor identity of
Lemma~\ref{lem:a} symbolically, that the coefficients of $(1+3x)/(1-x)^{4}$ agree with
$(e+1)(e+2)(4e+3)/6$ through $x^{24}$, and confirms
$\binom{e+3}{3}+3\binom{e+2}{3}=(e+1)(e+2)(4e+3)/6$ for $0\le e\le40$; and (iv) confirms
that the threefold transform of $4^{\omega}$ is multiplicative on all coprime pairs
$m,n<40$.

\begin{thebibliography}{9}
\bibitem{oeis} N.~J.~A. Sloane et al., \emph{The On-Line Encyclopedia of Integer
Sequences}, \url{https://oeis.org/A062368}, consulted 24 August 2026; also A001221 and
A007425.
\bibitem{apostol} T.~M. Apostol, \emph{Introduction to Analytic Number Theory}, Springer,
1976; Chapter 2 (Dirichlet multiplication) and Chapter 11 (Dirichlet series and Euler
products).
\end{thebibliography}
\end{document}
